\documentclass[11pt]{article}

\usepackage[margin=1.1in]{geometry}
\usepackage{amsmath,amssymb}
\usepackage{booktabs}
\usepackage{microtype}
\usepackage[hidelinks]{hyperref}
\usepackage{url}

\newcommand{\Ez}{$E_0$}

\title{When Structure Is Not Enough: A Preregistered Negative Evaluation\\
of a Training-Free Navigation Framework}
\author{Thomas Wehner\thanks{Developed and evaluated in a human--AI
collaboration; the AI systems involved are not authors. The collaboration
process, including its governance mechanisms, is described in
Section~\ref{sec:method}. Contact: repository issue tracker,
\url{https://github.com/Thomas66690815/E0-Framework}.}}
\date{August 2026}

\begin{document}
\maketitle

\begin{abstract}
We report the preregistered evaluation and closure of \Ez{}, a training-free
structural navigation framework built on two mechanisms: \emph{historization}
(persistent success/failure traces that modulate edge resistance) and a
\emph{geometric interference layer} (an exact discrete Helmholtz
decomposition inducing a $U(1)$ connection, path phases, and an
amplitude-based lookahead). Across 1{,}560 preregistered replicates on four
designed domain families at three scales under equal interaction budgets, we
find: (1)~the complete geometric stack contributes exactly $0.0$ over plain
historization --- a mechanism-level diagnosis shows the interference layer is
causally inactive at the action boundary; (2)~historization-only navigation
is not competitive with fairly trained standard baselines (success-adjusted
efficiency $0.208$ vs.\ $0.407$ for the per-instance median of tabular
Q-learning, UCB1, and random-restart greedy; paired difference $-0.199$,
95\,\% CI $[-0.206, -0.192]$); (3)~the single exception is a family of
wall-structured trap domains, where failure memory beats the baseline median
by $+0.260$ $[+0.248, +0.272]$. A separately preregistered calibration of
the framework's confidence-gated override selected \texttt{gate\_disabled}:
no threshold met its eligibility criteria. A final preregistered experiment
then tested the one surviving mechanism in its intended role --- persistent
tool-reliability memory for budget-constrained agents --- and passed:
$+48$--$54\,\%$ completed tasks over memoryless selection, with the honest
bound that a stateless sticky heuristic captures most of the stationary
value and the memory's real differentiator is recovery under drift. We
describe the preregistration, fair-baseline, causal-ablation, and
evidence-ledger methodology that produced these conclusions inside the same
repository that had strong incentives to avoid them, and we argue that the
mechanistic \emph{why} of the null result --- a conjunction of gating,
phase-regime degeneracy, and an implementation equivalence --- is more
informative than the null itself.
\end{abstract}

\section{Introduction}

\Ez{} began as an ambitious research program: derive learning, forgetting,
lookahead, and self-reflection from a single structural axiom, with no
probabilities, no training data, and no free parameters at the foundation.
Over 330 numbered commits it grew to ${\sim}36{,}000$ lines of production
code, 7{,}166 tests, and fourteen integrated layers.

This paper is not about that ambition. It is about what happened when the
program was forced to answer one question under preregistered rules:
\emph{does any of it beat the obvious alternatives under equal conditions?}

The honest answer is: mostly no, narrowly yes. We consider both halves of
that answer, and the machinery that extracted it, to be the program's actual
contribution. Negative results of this kind are rarely published with full
raw data, a frozen preregistration, causal ablations, and a mechanism-level
diagnosis of \emph{why} the flagship mechanism cannot act. This paper
provides all four.

Contributions:
\begin{enumerate}
  \item \textbf{A preregistered negative result} for a training-free
    structural navigation framework against fairly trained baselines, with
    all raw data retained (Section~\ref{sec:results2}).
  \item \textbf{A causal-ablation design} (five variants sharing one
    operationalization) that localizes all measured behavior in a single
    mechanism (Sections~\ref{sec:design} and~\ref{sec:results1}).
  \item \textbf{A mechanism diagnosis} explaining the exact conjunction that
    keeps the interference layer inert (Section~\ref{sec:why}).
  \item \textbf{A scoped positive finding}: edge-local failure memory wins
    where failures are persistent and structural
    (Section~\ref{sec:perfamily}).
  \item \textbf{A methodological account} of preregistration, evidence
    ledgers, and external audit inside a small human--AI research
    collaboration (Section~\ref{sec:method}).
  \item \textbf{An honestly bounded positive closure}: a final preregistered
    experiment on the surviving mechanism's intended use case, passed with
    its caveats made mandatory documentation (Section~\ref{sec:survives}).
\end{enumerate}

\section{The system under test}

\subsection{Historization}

\Ez{} represents a domain as a landscape of states and edges. Each edge
carries a difference $\Delta$ and a resistance $R$. The controller greedily
minimizes structural burden $S_{\mathrm{eff}} = \Delta \cdot
R_{\mathrm{eff}}$. After each transition the outcome is \emph{inscribed}:
successes ($U$) reduce effective resistance, failures ($F$) raise it, with
decay. This is the entire learning mechanism: no value function, no reward
propagation, no training phase. Credit assignment is strictly edge-local.

\subsection{The geometric layer}

The framework's distinctive construction: every flow field on a graph splits
uniquely (discrete Helmholtz) into a gradient part and a rotational part,
solved exactly via a graph-Laplacian system. The rotational component
induces a connection $\omega$, hence path phases $\Theta$, holonomy, and
curvature. A lookahead layer sums complex amplitudes $\Psi = e^{-S} \cdot
e^{i\Theta}$ over forward path families; interference is meant to
concentrate support on goal-reaching actions. A confidence-gated override
lets this layer replace the greedy choice when its margin is large.

To our knowledge the decomposition-induced phase is novel in this context
(quantum walks and Projective Simulation interfere amplitudes on graphs but
do not derive the phase from an orthogonal field decomposition). Novelty, as
Section~\ref{sec:results2} shows, is not value.

\section{Preregistered design}\label{sec:design}

The full protocol (\texttt{E0-G1-v1}) was frozen before any full-budget run,
including domains, splits, budgets, ablations, metrics, statistics, and
error rules. Key elements:

\begin{itemize}
  \item \textbf{Domain families (4):} \texttt{wall\_grid} (walls, dead ends,
    traps), \texttt{trap\_grid\_v2} (dense trap structure),
    \texttt{decoy\_dag} (attractive decoy paths),
    \texttt{nonstationary\_parallel} (mid-run executor switches).
  \item \textbf{Scales:} $N \in \{100, 500, 1000\}$ states; exact-$N$
    generators.
  \item \textbf{Splits:} development seeds 0--9 (implementation checks,
    control selection); holdout seeds 1000--1029, never opened.
  \item \textbf{Budgets:} identical interaction and episode budgets for
    every method; 10 adaptation $+$ 20 evaluation episodes per replicate;
    hard per-episode (60\,s) and per-replicate (1{,}800\,s) deadlines in
    killable subprocesses.
  \item \textbf{Causal ablations (5):} \texttt{A\_HIST} (historization
    only), \texttt{B\_INCOHERENT} (amplitudes without phase coherence),
    \texttt{C\_THETA\_ZERO} (coherent, phase forced to zero),
    \texttt{D\_U1\_PHASE} ($U(1)$ phase active), \texttt{E\_FULL\_GEOMETRY}
    (full stack). All five share one frozen operationalization, path family,
    and override rule, so any outcome difference is attributable to the
    ablated component.
  \item \textbf{Baselines (8, fairly configured):} tabular Q-learning,
    UCB1-edge, random-restart greedy (the three \emph{competitive}
    comparators, combined per instance as their median), plus
    $\varepsilon$-greedy, memoryless greedy, uniform random, and
    map-informed A*/D*-Lite as upper references excluded from the
    comparator.
  \item \textbf{Statistics:} per-instance pairing, stratified cluster
    bootstrap (10{,}000 resamples), fixed seeds, preregistered family-level
    criteria and aggregation rules.
\end{itemize}

Execution: GitHub Actions, 240 deterministic batches, atomic per-replicate
shards, fail-closed consolidation; 1{,}560/1{,}560 replicates completed,
zero infrastructure failures, three valid algorithmic timeouts (all
\texttt{E\_FULL\_GEOMETRY}). The consolidated bundle (raw runs, 31{,}200
evaluation episodes, frozen configs, environment, manifest with SHA-256 file
digests) is retained in-repository.

\section{Results I: the geometry is causally inert}\label{sec:results1}

Every geometry variant scores identically to plain historization
(Table~\ref{tab:geometry}).

\begin{table}[ht]
\centering
\begin{tabular}{lrr}
\toprule
Comparison & Mean difference & 95\,\% CI \\
\midrule
\texttt{E\_FULL\_GEOMETRY} $-$ \texttt{A\_HIST} (overall, every family) & $0.000$ & $[0.000, 0.000]$ \\
\texttt{D\_U1\_PHASE} $-$ \texttt{C\_THETA\_ZERO} (phase attribution) & $0.000$ & $[0.000, 0.000]$ \\
\bottomrule
\end{tabular}
\caption{Causal ablation outcomes; paired stratified bootstrap.}
\label{tab:geometry}
\end{table}

This is not approximate equality; across all 17{,}929{,}640 persisted
evaluation decisions of the B--E variants there is not a single executed
override, not a single interfering or wrapped phase regime, and not a single
divergent selected action. The only measured difference is cost: the
full-geometry variants spend 139--146\,s median wall-time per replicate
against 8.8\,s for \texttt{A\_HIST} --- a ${\sim}16\times$ compute overhead
for byte-identical outcomes.

\section{Results II: historization alone is not competitive}\label{sec:results2}

\subsection{Overall}

Success-adjusted efficiency (primary metric), development split, 120 paired
instances (Table~\ref{tab:methods}).

\begin{table}[ht]
\centering
\begin{tabular}{lrrr}
\toprule
Method & Efficiency & Goal rate & Median wall-time \\
\midrule
A*/D*-Lite (map-informed reference) & 0.875 & 0.875 & 0.2--0.3\,s \\
Random-restart greedy & 0.464 & 0.538 & 3.5\,s \\
UCB1-edge & 0.393 & 0.396 & 3.2\,s \\
Q-learning (tabular) & 0.385 & 0.771 & 4.3\,s \\
Uniform random & 0.273 & 0.349 & 4.3\,s \\
\textbf{A\_HIST (historization only)} & \textbf{0.208} & \textbf{0.208} & 8.8\,s \\
$\varepsilon$-greedy & 0.158 & 0.158 & 6.7\,s \\
Memoryless greedy & 0.125 & 0.125 & 7.2\,s \\
\bottomrule
\end{tabular}
\caption{Method aggregates over the full development matrix.}
\label{tab:methods}
\end{table}

\texttt{A\_HIST} vs.\ the competitive baseline median: $-0.199$
$[-0.206, -0.192]$, goal rate $-27.8$ percentage points. That uniform random
outperforms the framework overall is the kind of sentence a project only
writes down if its rules force it to; our rules forced it, and the number is
real. (It reflects the low goal rates of edge-local methods on decoy- and
nonstationary-dominated families, where random restarts escape what
inscribed penalties do not.)

\subsection{Per family: the niche is real}\label{sec:perfamily}

\begin{table}[ht]
\centering
\begin{tabular}{lrrl}
\toprule
Family & \texttt{A\_HIST} $-$ median & 95\,\% CI & Verdict \\
\midrule
\texttt{decoy\_dag} & $-0.763$ & $[-0.780, -0.744]$ & decisive loss \\
\texttt{nonstationary\_parallel} & $-0.283$ & $[-0.295, -0.270]$ & decisive loss \\
\texttt{trap\_grid\_v2} & $-0.008$ & $[-0.021, 0.000]$ & both fail (${\sim}0$) \\
\texttt{wall\_grid} & $\mathbf{+0.260}$ & $[+0.248, +0.272]$ & \textbf{clear win} \\
\bottomrule
\end{tabular}
\caption{Family-level comparison against the fair baseline median.}
\label{tab:family}
\end{table}

One of four families meets the preregistered competitiveness bar; the gate
required three. But the \texttt{wall\_grid} result is a clean, replicable
positive: where failure is \emph{persistent and structural} --- walls and
dead ends that stay walls and dead ends --- cheap edge-local failure memory
beats trained baselines that must rediscover the same dead ends across
episodes. Where failure is \emph{misleading} (decoys) or \emph{transient}
(nonstationarity), inscribed penalties are at best useless and at worst
poison.

\subsection{The override gate calibrates to ``off''}\label{sec:gate}

Independently of the gate experiment, the confidence-gated override was
given its own preregistered calibration (2{,}880 tasks, 12 candidate
policies, clustered bootstrap, Holm-corrected eligibility across 77
constraints). Outcome: no active candidate was eligible; the preregistered
fail-closed selection was \texttt{gate\_disabled}. A caveat we discovered
afterwards and report openly: 73\,\% of calibration replicates hit a timeout
whose label conflated parent computation with diagnostic branch
instrumentation, so the calibration cannot cleanly rank thresholds; it can
only support ``no threshold met the criteria as measured.''

\subsection{Confirmed structural limits}

Two limits established earlier by falsification benchmarks stand: dense
branching ($b \geq 3$) defeats the penalty mechanism combinatorially (F3),
and non-Markov dependencies cannot be learned because credit assignment is
edge-local (F4) --- the system avoids the trap but cannot acquire the
required sequence.

\section{Why the interference layer cannot act}\label{sec:why}

The neutral result in Section~\ref{sec:results1} is not a tie between active
mechanisms; it is a conjunction of three blockers, each independently
sufficient on these domains:

\begin{enumerate}
  \item \textbf{The gate never opens.} The preregistered override threshold
    is $0.85$; the maximum conflict confidence ever observed when lookahead
    disagreed with greedy was $0.42$ (probe over 9{,}600 decisions).
    High-confidence states exist, but only when lookahead already
    \emph{agrees} with greedy --- precisely when an override changes
    nothing. The threshold lies outside the support of the conflict
    distribution.
  \item \textbf{The phase never leaves the gradient regime.} All
    4{,}482{,}410 evaluation decisions per variant sit in the gradient
    regime; circulation on these families is too weak to produce interfering
    or wrapped phases. No circulation $\Rightarrow$ no phase $\Rightarrow$
    interference degenerates to plain summation (consistent with an earlier
    finding that at typical weights interference is purely constructive;
    destructive cancellation needs phase gaps ${\approx}\,\pi$, outside the
    stable ranking range).
  \item \textbf{An implementation equivalence.} For the current
    influence-map construction, the packaged E path is mathematically
    equivalent to D, so even a phase-active regime would not separate the
    top two variants.
\end{enumerate}

The general lesson: a mechanism can be mathematically real, correctly
implemented, fully tested --- and still \emph{causally unreachable} under
the system's own operating conditions. Only a causal ablation run at full
budget, not unit tests and not demos, exposes this.

\section{Method: making a repository falsify itself}\label{sec:method}

The interesting engineering problem was not statistical but institutional:
the same collaboration that spent months building the geometry had to be
able to kill it. Mechanisms that made this work:

\begin{itemize}
  \item \textbf{Preregistration with freezes.} Domains, budgets, criteria,
    statistics, and seeds frozen before full-budget execution; post-freeze
    edits require labeled errata (one occurred, and an external audit caught
    its missing label).
  \item \textbf{Fair baselines with an audit trail.} Every baseline
    configuration frozen and audited; map-informed methods explicitly
    excluded from the comparator; the embarrassing baselines (uniform
    random, memoryless greedy) included.
  \item \textbf{An evidence ledger.} Every strategic claim maps to a ledger
    entry with status, sources, limitations, and raw artifacts; claims are
    downgraded in place when contradicted.
  \item \textbf{Split discipline.} Development seeds for selection, a
    holdout that was never opened --- and, at closure, the honest refusal to
    open it merely to ceremonialize a conclusion the development data had
    already fixed.
  \item \textbf{Hard execution boundaries.} Killable per-episode and
    per-replicate deadlines after a first pilot showed unbounded timeouts
    silently corrupting attribution; the aborted pilot was quarantined as
    engineering-only evidence.
  \item \textbf{External audit.} An independent review verified hash chains,
    freeze integrity, count consistency, and authorization boundaries, and
    surfaced the one unreported result in the pipeline (the baseline
    comparison in Section~\ref{sec:results2}), which was then retained and
    reported before its CI artifact expired.
\end{itemize}

We also name the anti-pattern we had to resist: \textbf{infrastructure
regress} --- answering every ambiguous or unwelcome measurement by building
a more rigorous measurement apparatus instead of drawing the available
conclusion. The override-gate line (Section~\ref{sec:gate}) consumed sixteen
commits of exemplary methodology to calibrate a knob on a mechanism that
Section~\ref{sec:results1} shows never fires. Rigor was necessary; more
rigor was not the missing ingredient --- a decision was.

\section{Limitations}

\begin{itemize}
  \item All results are on designed synthetic families; the strong
    \texttt{wall\_grid} claim has not been validated on an external domain.
  \item The gate was closed on development evidence; no holdout PASS/FAIL
    exists. We report this as a strategic closure, not a formal gate
    outcome.
  \item Single implementation, single research group; no independent
    replication. The full raw data, frozen configs, and environment
    manifests are retained to make replication cheap.
  \item The baseline set is standard but not exhaustive; no deep-RL or
    LLM-based navigation baselines were included (they violate the
    equal-budget, training-free comparison frame).
  \item The final experiment (Section~\ref{sec:survives}) models tool
    failure as seeded Bernoulli processes; no production tool ecosystem was
    measured, and no LLM was in the loop by design.
\end{itemize}

\section{What survives: the final experiment}\label{sec:survives}

A ${\sim}600$-line, dependency-free failure-memory library
(\texttt{reliability\_memory}) embodies the one mechanism that carried all
measured behavior. Its remaining product hypothesis --- that an agent which
persistently remembers \emph{which tools and paths fail} outperforms the
same agent without memory, in settings where no training budget exists ---
is deliberately \emph{not} addressed by the experiments above, whose
baselines received training budgets. We closed the program with one
preregistered decision experiment on this hypothesis (protocol
\texttt{E0-WP6-RELMEM-v1}; design frozen and externally approved before
implementation; shipped defaults, no tuning; the library bound by SHA-256).

Setup: simulated tool ecosystems --- tasks of five steps, four redundant
tools per step, 1{,}000 calls per replicate --- in three regimes: persistent
failure structure (R1), mid-run drift (R2), and task-type-dependent
reliability (R3). Four arms: the memory store, memoryless uniform selection,
a stateless \emph{sticky} heuristic (``keep the tool that last worked''),
and an oracle upper reference; 30 paired seeds per regime,
10{,}000-resample paired bootstrap. Losing to the sticky one-liner in R1 was
preregistered as failure.

\begin{table}[ht]
\centering
\begin{tabular}{lrrrr}
\toprule
Regime & Memory & No memory & Sticky & Oracle \\
\midrule
R1 persistent & \textbf{125.9} & 85.1 & 122.8 & 153.3 \\
R2 drift & \textbf{124.2} & 80.9 & 118.3 & 148.0 \\
R3 context & 85.0 & 76.9 & \textbf{105.2} & 139.3 \\
\bottomrule
\end{tabular}
\caption{Final experiment: completed tasks per 1{,}000 tool calls.}
\label{tab:wp6}
\end{table}

The verdict under the frozen criteria is \textbf{PASS}: $+48.0\,\%$ over
memoryless selection in R1 (95\,\% CI of the absolute difference
$[+35.3, +46.8]$), $+53.5\,\%$ in R2, $+10.5\,\%$ in R3, and a positive CI
against sticky in R1 ($+3.2$ $[+1.4, +4.9]$). The honest bounds are part of
the result: in stationary regimes the sticky heuristic captures most of the
achievable value (memory's edge is only $+2.6\,\%$, a mandatory caveat now
permanent in the library's documentation); the memory's real differentiator
is \emph{drift}, where decaying traces recover and a last-success pointer
does not ($+5.0\,\%$ over sticky, CI $[+3.9, +8.0]$); and where reliability
depends on task context, a context-keyed sticky heuristic beat the
deliberately context-free store by $19.2\,\%$ --- the state key must carry
the context. The program therefore ends with a small maintained tool whose
claims are exactly as large as its evidence, and no larger.

\section{Conclusion}

We set out to derive intelligent navigation from structure alone and built
the apparatus to test whether we had. We had not: the geometric layer is
causally inert on its own benchmark, and the memory mechanism alone loses to
thirty-year-old baselines everywhere except the one environment class whose
failure structure matches its assumptions. The same apparatus, pointed at
the surviving mechanism's actual use case, returned a scoped positive with
its caveats attached. We consider the preregistered negative, its
mechanistic explanation, the honestly bounded positive, and the
demonstration that a small research program can be made to falsify itself,
to be worth more than the framework was.

\section*{Data availability}

All raw runs, episode records, frozen configurations, environment manifests,
and SHA-256 manifests are public:
\texttt{artifacts/g1/E0-G1-v1/development/run\_30526724307/} (gate
experiment) and \texttt{artifacts/wp6/E0-WP6-RELMEM-v1/} (final experiment)
in \url{https://github.com/Thomas66690815/E0-Framework}. Evidence ledger:
\texttt{docs/E0\_EVIDENCE\_LEDGER\_v1.json}. Closure record:
\texttt{docs/E0\_G1\_CLOSURE\_v1.md}.

\end{document}
